\section{The keystone: the recurring four is one object}

This section records claim nodes and dependency edges.

\begin{proposition}[Bare bijection]\label{thm:bij}\lean{Infomaterial.Bij}\leanok\uses{def:articulating}
A \texttt{Bij}~$\alpha$~$\beta$ is a function $\mathrm{to} : \alpha \to \beta$ together with an inverse $\mathrm{inv} : \beta \to \alpha$ and the two round-trip laws $\mathrm{inv}\,(\mathrm{to}\,a) = a$ and $\mathrm{to}\,(\mathrm{inv}\,b) = b$. It is the dependency-free witness of ``these two finite objects are the same,'' used everywhere the chapter asserts an identity of readings.
\end{proposition}
\begin{proposition}[Articulating.Corner as status pair]\label{thm:corner-status}\lean{Infomaterial.cornerStatus}\leanok\uses{thm:bij, def:articulating, thm:four-bilattice}
The bijection $\texttt{cornerStatus} : \texttt{Bij}\;\texttt{Articulating.Corner}\;(\texttt{Bool} \times \texttt{Bool})$ sends an articulation phase to its status pair $(\text{token-articulated}, \text{claim-articulated}) = (\text{positive}, \text{negative})$ and back. The four articulation phases are exactly the four two-sided status pairs; the identification is proven by exhausting the four corners.
\end{proposition}
\begin{proposition}[The four predicate generators]\label{thm:generator}\lean{Infomaterial.Generator}\leanok\uses{thm:predication, thm:predication-converse}
\texttt{Generator} is the four-constructor inductive of morphism classes of the predication category: specialization, converse, subject-projection, object-projection. Its enumeration \texttt{Generator.all} has length four (\texttt{Generator.card\_eq\_four}); these are the recurring four met as the generators of predicate morphisms.
\end{proposition}
\begin{theorem}[The four-fold identity]\label{thm:four-fold}\lean{Infomaterial.four_fold}\leanok\uses{thm:corner-status, thm:generator}
\texttt{four\_fold} states that \texttt{Articulating.Corner.all} and \texttt{Generator.all} each have length four: every incarnation of the recurring four has exactly four elements, with no fifth reading.
\end{theorem}
\begin{theorem}[Charge is phase, definitionally]\label{thm:charge-is-phase}\lean{Infomaterial.charge_is_phase}\leanok\uses{thm:four-bilattice, def:articulating}
\texttt{charge\_is\_phase} proves $(c : \texttt{Charge}) = c$ for every corner $c$ --- because \texttt{Charge} is \emph{defined} to be \texttt{Articulating.Corner}. A bilattice charge is not merely isomorphic to an articulation phase; it \emph{is} one, re-met as a value rather than as a stage of descent. This definitional collapse is the spine of the keystone.
\end{theorem}
\begin{proposition}[The canonical four-corner object]\label{thm:four-object}\lean{Infomaterial.Articulating.FourCorner}\leanok\uses{thm:corner-status, thm:charge-is-phase}
\texttt{Articulating.FourCorner} is the abbreviation $\texttt{Bool} \times \texttt{Bool}$: two binary slot-presences. This is the headline object of the chapter: the canonical carrier of the quartet pattern. The bijection \texttt{Articulating.FourCorner.cornerBij} identifies it with \texttt{Articulating.Corner}, and since \texttt{Charge}~$=$~\texttt{Articulating.Corner}, also with \texttt{Charge}: the canonical four-corner \emph{is} the bilattice value object.
\end{proposition}
\begin{proposition}[A cell's four-corner]\label{thm:cell-fourcorner}\lean{Infomaterial.Articulating.Cell.fourCorner}\leanok\uses{thm:four-object, def:articulating}
\texttt{Articulating.Cell.fourCorner} reads the articulation phase of a cell as the four-corner $(\text{tokenArticulated}, \text{claimArticulated})$. It is the bridge across which a stage of the descent becomes a value of the canonical object.
\end{proposition}
\begin{theorem}[Conflation: articulating $\cong$ positing]\label{thm:articulation-phase-is-charge}\lean{Infomaterial.articulation_phase_is_charge}\leanok\uses{thm:cell-fourcorner, thm:four-object}
\texttt{articulation\_phase\_is\_charge} proves that \texttt{Articulating.FourCorner.toCorner} of a cell's four-corner equals the cell's corner --- and since \texttt{Articulating.Corner}~$=$~\texttt{Charge}, equals its bilattice charge. The two-sided $(\text{token}, \text{claim})$ partition that articulating produces and the four-point value that positing admits are one object. The two verbs are the same four-corner cell met at two depths --- a theorem, not a comment.
\end{theorem}
\begin{theorem}[The descent's start is one point]\label{thm:potential-is-kbot}\lean{Infomaterial.potential_is_kBot}\leanok\uses{thm:articulation-phase-is-charge, thm:settling}
\texttt{potential\_is\_kBot} proves that the initial (potential) cell of any utterance sits at the $(\texttt{false}, \texttt{false})$ corner, which is the knowledge-bottom charge $\bot_k$ that positing seeds. The start of the descent is a single point in both the articulation reading and the value reading.
\end{theorem}
